
\section{z-Transformation / Zeitdiskrete Übertragungsfunktion}

\subsection{Folgen (und deren z-Transformation)}{230}

Eine zeitliche Folge $f(k)$ bricht typischerweise nicht ab und wird folgendermassen beschrieben:

\vspace{-0.2cm}

$$ f(k) = \{1, \, a, \, a^2, \, a^3, \, \ldots \} = a^k \qquad \text{für } k \geq 0 \quad k \text{ ganzzahlig} $$


\subsubsection{Spezielle Folgen: Abbrechende Folgen}

Bei einer abbrechenden Folge liegen \textbf{alle Pole bei Null}. \\
Man spricht in diesem Fall von einem \textbf{FIR-Filter} (Finite Response Filter)


\example{Abbrechende Folge und deren z-Transformierte}

\vspace{-0.2cm}

$$ f(k) = \left\{ f(0), \, f(1), \, f(2), \, f(3), \, 0, \, 0, \, 0, \, \ldots  \right\} $$

\vspace{-0.2cm}

$$ F(z) = f(0) + \frac{f(1)}{z} + \frac{f(2)}{z^2} + \frac{f(3)}{z^3} = \frac{f(0) \, z^3 + f(1) \, z^2 + f(2) \, z + f(3)}{z^3}$$



\subsection{z-Transformation -- Definition}{232}

Die Überführung eine Folge in den Frequenzbereich erfolgt mittels z-Tranaformation und ist definiert als

\vspace{-0.2cm}

$$ \boxed{ F(z) = \mathcal{Z}\{f(k)\} = \sum\limits_{k=0}^{\infty} f(k) \cdot z^{-k} \quad \text{mit der komplexen Variablen } z } $$


\textrightarrow\ Häufig werden \textbf{Tabellen} und die \textbf{Eigenschaften} der z-Transformation verwendet \\
\textrightarrow\ Siehe \ref{zTransformierte einger Funktionen} und \ref{Eigenschaften der z-Transformation}



\subsection{z-Transformierte einiger Funktionen}{232}
\label{zTransformierte einger Funktionen}

Alle Folgen $f(k)$ sind nur für für $k \geq 0$ und ganzzahlige $k$ definiert.

\smallskip

\begin{minipage}[t]{0.44\columnwidth}
    \begin{center}
        \renewcommand{\arraystretch}{1.4}
        \begin{tabular}{c c c}
            \toprule
            $\bm{f(k)}$                             & $\laplace$    & $\bm{F(z)}$                                       \\
            \midrule
            $\delta(k)$                             & $\laplace$    & $1$                                               \\
            \midrule
            $\varepsilon(k)$                        & $\laplace$    & $\frac{z}{z-1}$                                   \\
            \midrule
            $k$                                     & $\laplace$    & $\frac{z}{(z-1)^2}$                               \\
            \midrule            
            $k^2$                                   & $\laplace$    & $\frac{z(z+1)}{(z-1)^3}$                          \\
            \midrule
            $a^k$                                   & $\laplace$    & $\frac{z}{z-a}$                                   \\
            \bottomrule
        \end{tabular}
        \renewcommand{\arraystretch}{1}
    \end{center}
\end{minipage}
\hfill
\begin{minipage}[t]{0.54\columnwidth}
    \begin{center}
        \renewcommand{\arraystretch}{1.4}
        \begin{tabular}{c c c}
            \toprule
            $\bm{f(k)}$                             & $\laplace$    & $\bm{F(z)}$                                           \\
            \midrule
            $\cos(bk)$                              & $\laplace$    & $\frac{z(z - \cos(b))}{z^2 - 2z \cos(b) + 1}$         \\
            \midrule
            $\sin(bk)$                              & $\laplace$    & $\frac{z \sin(b)}{z^2 - 2z \cos(b) + 1}$              \\
            \midrule
            $k \cdot a^k$                           & $\laplace$    & $\frac{az}{(z-a)^2}$                                  \\
            \midrule
            $k \cdot \cos(bk)$                      & $\laplace$    & $\frac{z(z - a \,\cos(b))}{z^2 - 2az \cos(b) + a^2}$  \\
            \midrule
            $k \cdot \sin(bk)$                      & $\laplace$    & $\frac{z a \sin(b)}{z^2 - 2az \cos(b) + a^2}$         \\
            \bottomrule
        \end{tabular}
        \renewcommand{\arraystretch}{1}
    \end{center}
\end{minipage}

\smallskip

\textbf{Hinweis:} $\varepsilon(k)$ ist eine Folge von $1$ern \\
\textbf{Hinweis:} Oft wird $\frac{F(z)}{z}$ in der Tabelle gefunden

\para{DC-Verstärkung}

\[
G_{DC}
=
G(z)\Big|_{z=1}
\]

\subsection{Eigenschaften der z-Transformation}{235}
\label{Eigenschaften der z-Transformation}

\resizebox{\columnwidth}{!}{
    \begin{tabular}{ccc}
        \toprule
        \textbf{Zeitbereich}                                                        & \textbf{Bildbereich}                          & \textbf{Operation im Zeitbereich}     \\
        \midrule
        $a_1 f_1(k) + a_2 f_2(k)$                                                   & $a_1 F_1(z) + a_2 F_2(z)$                     & Superposition                         \\
        \midrule
        $f(k+1)$                                                                    & $z F(z) - z \, f(0)$                          & Linksverschiebung...                  \\
        \midrule
        $f(k+2)$                                                                    & $z^2 F(z) - z^2 f(0) - z f(1)$          & ...mehrfach angewendet                \\
        \midrule
        $f(k-m) \quad m \geq 0$                                                     & $\frac{1}{z^m} F(z)$                          & Rechtsverschiebung                    \\
        \midrule
        $\sum\limits_{\nu=0}^k f(\nu)$                                              & $\frac{z}{z-1} F(z)$                          & Partialsummen                         \\
        \midrule
        $\sum\limits_{\nu=0}^k f_1(\nu) \cdot f_2(k - \nu)$                         & $F_1(z) \cdot F_2(z)$                         & Faltung                               \\
        \bottomrule
    \end{tabular}
}


\subsubsection{Grenzwertsätze der z-Transformation}{235}

\textbf{Achtung:} Grenzwertsätze dürfen nur angewendet werden, wenn Grenzwerte existieren! \\
\textrightarrow\ In der Regelungstechnik eigentlich immer gegeben...

\begin{center}
    \begin{tabular}{l c ll}
        \toprule
        \textbf{Zeitbereich}                &       & \textbf{Bildbereich}                                      & \textbf{Bezeichnung}      \\
        \midrule
        $f(0)$                              & $=$   & $\lim\limits_{z \to \infty} F(z)$                         & Anfangswerttheorem        \\
        \midrule
        $\lim\limits_{k \to \infty} f(k)$   & $=$   & $\lim\limits_{z \to 1} \left[ (z-1) \cdot F(z) \right]$   & Endwerttheorem            \\
        \bottomrule
    \end{tabular}
\end{center}


\subsection{Rücktransformation}
\label{Rücktransformation}

Die Rücktransformation kann auf mehrere Arten erreicht werden:

\smallskip

\begin{enumerate}
    \item Direktes benutzen von Tabellen in Rückwärtsrichtung
    \item Zerlegung von $F(z)$ in einfachere, tabellierte Ausdrücke \textrightarrow\ Partialbruchzerlegung
    \item Polynomdivision bzw. Reihenentwicklung von $f(z)$ in negativen Potenzen von $z$
    \item (Allgemeine Umkehroperation zur Definition der z-Transformation)
\end{enumerate}

\smallskip 

\textbf{Bemerkung zu 2.:} Statt $F(z)$ zu zerlegen wird der Ausdruck $\frac{F(z)}{z} = s_1 + s_2 + \ldots$ zerlegt.
Anschliessend werden die Ausdrücke $z \cdot s_i$ rücktransformiert, deren Summe $f(k)$ ergibt. \\
Somit ist die Wahrscheinlicheit höher, den benötigten Term in der Tabelle zu finden. 



\subsection{Beispiele zur z-Transformation}

\example{Geometrische Folge im Frequenzbereich}

Die geometrische Folge ist das Pendant zur zeitkontinuierlichen Exponentialfunktion.

\vspace{-0.2cm}

\[
    F(z)    = \sum\limits_{k=0}^{\infty} a^k \cdot z^{-k} = \sum\limits_{k=0}^{\infty} a^k \cdot \left( \frac{a}{z} \right)^k
            = 1 + \frac{a}{z} + \left( \frac{a}{z} \right)^2 + \left( \frac{a}{z} \right)^3 + \ldots 
\]

Durch eine geeignete Umformung kann die Folge anders dargestellt werden und anschliessend nach $F(z)$ aufgelöst werden.

\vspace{-0.2cm}

\[
    F(z) - 1 = \frac{a}{z} \cdot F(z) \quad \Leftrightarrow \quad 
    1 = F(z) \left( 1 - \frac{a}{z} \right) \quad \Leftrightarrow \quad 
    F(z) = \frac{z}{z-a}
\]


\example{Periodische Folge}

\vspace{-0.2cm}

$$ f_5(k) = \{ \underbrace{0, \, 0, \, 0, \, 0, \,}_{n \text{ Glieder}} \underbrace{ h(0), \, h(1), \, \ldots, \, h(m-1)}_{m \text{ Glieder}}, \underbrace{ h(0), \, h(1), \, \ldots, \, h(m-1)}_{m \text{ Glieder}}, \, \ldots \} $$

$$ F_5(z) = z^{-m} \cdot F_5(z) + z^{-n} \cdot H(z)  \qquad \Rightarrow\ \qquad F_5(z) = \frac{1}{z^n} \frac{z^m}{z^m -1} H(z) $$


\textrightarrow\ $H(z)$ entspricht der z-Transformierten von $\left\{  h(0), \, h(1), \, \ldots, \, h(m-1) 
, \, 0, \, 0, \, \ldots \right\}$


\example{Ausgangssignal bestimmen}

Der Ausgang entspricht der \textbf{Faltung} vom Eingangssignal und Impulsantwort: 
$$ y(k) = u(k) * g(k) $$
Bei jedem Schritt $k$ wird die Impulsantwort mit dem aktuellen $u(k)$ 'skaliert' um $k$ Schritte nach rechts geschoben.

\smallskip

Für $u(k) = \{ 1, \, 0, \, 1, \, 0 , \, 0, \, \ldots \}$ und $g(k) = \{ 6, \, -2, \, -2, \, -2 , \, 0, \, 0, \, \ldots \}$ soll das Ausgangssignal $y(k)$ berechnet werden.

\begin{ctabular}{r l c c c c c c c}
    $u(0) \cdot g(k) =$     & $\{ 6,$   & $-2,$ & $-2,$ & $-2,$ & $0,$  & $0,$  & $0,$  & $\ldots \}$   \\
    $u(1) \cdot g(k-1) =$   & $\{ 0,$   & $0,$  & $0,$  & $0,$  & $0,$  & $0,$  & $0,$  & $\ldots \}$   \\
    $u(2) \cdot g(k-2) =$   & $\{ 0,$   & $0,$  & $6,$  & $-2,$ & $-2,$ & $-2,$ & $0,$  & $\ldots \}$   \\
    $y(k) =$                & $\{ 6,$   & $-2,$ & $4,$  & $-4,$ & $-2,$ & $-2,$ & $0,$  & $\ldots \}$   \\
\end{ctabular}



\example{Lösen einer linearen Differenzengleichung im Bildbereich}

Wie eine Differentialgleichung mittels Laplace-Transformation wird auch eine Differenzengleichung mittels z-Transformation in drei Schritten gelöst:

\smallskip

\begin{enumerate}
    \item Überführung der Differenzengleichung (DG) in den Bildbereich (\textrightarrow\ z-Transformation)
    \item Lösen des resultierenden algebraischen Ausdrucks
    \item Rücktransformation von $Y(z)$ in den Zeitbereich (\textrightarrow\ siehe Abschnitt \ref{Rücktransformation})
\end{enumerate}


$$ \text{Differenzengleichung:} \qquad  a_1 \cdot y(k+1) + a_0 \cdot y(k) = b_0 \cdot u(k) $$

\vspace{-0.2cm}

$$ \bm{1.} \quad a_1 \cdot z \cdot Y(z) - a_1 \cdot z \cdot y(0) + a_0 \cdot Y(z) = b_0 \cdot U(z) $$
$$ \bm{2.} \quad Y(z) = \frac{b_0}{a_0 + a_1 \cdot z} \cdot U(z) + \frac{z}{a_0 + a_1 \cdot z} \cdot y(0) \quad \underrightarrow{\bm{3.} \text{ Tabellen}} \quad y(z) $$ 


\subsection{Zeitdiskrete Übertragungsfunktion}{238-239}

Die \textbf{zeitdiskrete Übertragungsfunktion} $\bm{G(z)}$ beschreibt das Verhalten zwischen Ausgang und Eingang eines diskreten 
Systems und ergibt sich aus der Differenzengleichung im Zeitbereich.

\vspace{-0.4cm}

$$ G(z) = \frac{Y(z)}{U(z)} = \frac{b_m z^m + b_{m-1} z^{m-1} + \ldots + b_0}{z^n + a_{n-1} z^{n-1} + \ldots + a_0} = K \cdot \frac{(z - z_1) (z - z_2) \ldots (z - z_m)}{(z - p_1) (z - p_2) \ldots (z - p_n)} \quad (m \leq n)$$

\textbf{Hinweis:} Für $m > n$ wäre $G(z)$ \textbf{akausal} und somit nicht realisierbar

\smallskip

\textrightarrow\ Gerechnet werden kann mit $G(z)$ gleich wie mit $G(s)$


\subsubsection{DC-Gain}

Der DC-Gain ist leicht anders definiert als bei der Laplace-Transformation

\smallskip

\begin{minipage}[t]{0.48\columnwidth}
    \paragraph{z-Transformation}

    \vspace{-0.2cm}

    $$ \text{DC-Gain} = G(z) \Big|_{z=1} $$
\end{minipage}
\hfill
\begin{minipage}[t]{0.48\columnwidth}
    \paragraph{Laplace-Transformation}

    \vspace{-0.2cm}

    $$ \text{DC-Gain} = G(s) \Big|_{s=0} $$
\end{minipage}



\subsection{Stabilität einer zeitdiskreten Übertragungsfunktion}{240-241}

Für die Untersuchung der Stabilität eines Systems wird der \textbf{Nenner $\bm{N(z)}$ der Übertragungsfunktion} $\bm{G(z)}$ betrachtet.

\smallskip

% SigSys2 mit Anpassungen
Für \textbf{Grenzstabilität} gilt eine \textbf{UND-Verknüpfung} der aufgeführten Punkte. \\
Für \textbf{Stabilität und Instabilität} gilt eine \textbf{ODER-Verknüpfung} der aufgeführten Punkte.

\vspace{0.1cm}

\begin{outline}
    \1 \textbf{Stabil:} 
        \2 Alle Polstellen innerhalb des Einheitskreises \textrightarrow\ $\abs{z} < 1$
        \2 Keine Polstellen vorhanden 
    \1 \textbf{Asymptotisch stabil:}
        \2 Alle Polstellen innerhalb des Einheitskreises \textrightarrow\ $\abs{z} < 1$
    \1 \textbf{Grenzstabil:} 
        \2 \textbf{Keine} Polstellen ausserhalb des Einheitskreises \textrightarrow\ $\abs{z} > 1$
        \2 Mindestens eine \textbf{einfache Polstelle} auf dem Einheitskreis \textrightarrow\ $\abs{z} = 1$
        \2 \textbf{Keine doppelten} Polstellen auf dem Einheitskreis \textrightarrow\ $\abs{z} = 1$
    \1 \textbf{Instabil:} 
        \2 Mindestens eine Polstelle ausserhalb des Einheitskreises \textrightarrow\ $\abs{z} > 1$
        \2 Mindestens eine \textbf{mehrfache Polstelle} auf dem Einheitskreis \textrightarrow\ $\abs{z} = 1$
\end{outline}

\smallskip

\textbf{Hinweis:} 'Stabil' ist ein Synonym für 'asymptotisch stabil' \\
\textbf{Hinweis:} Tustin erhält die Stabilitätseigenschaften eines Systems



\subsubsection{Jury-Kriterium}{241}

Als Pendant zum Routh-Kriterium im Laplace-Bereich existiert in der z-Ebene das Jury-Kriterium.
Somit kann vereinfacht festgestellt werden, ob sich alle Pole des Polynoms $N(z)$ \textbf{innerhalb des Einheitskreises} befinden und das System somit \textbf{stabil} ist.

\renewcommand{\arraystretch}{1.6}
\begin{ctabular}{|c|c c|}
    \hline
    \textbf{Systemordnung}      & \multicolumn{2}{c|}{\textbf{Stabilitätsbedingungen für $\bm{N(z) = z^n + a_{n-1}z^{n-1} + \cdots + a_0$}}} \\
    \hline
    $n = 1$ & & $|a_0| < 1$ \\
    \hline
    $n = 2$ & 
    $\begin{aligned}
    a_0 + a_1 + 1 &> 0 \\
    a_0 - a_1 + 1 &> 0
    \end{aligned}$ & $|a_0| < 1$ \\
    \hline
    $n = 3$ & 
    $\begin{aligned}
    a_0 + a_1 + a_2 + 1 &> 0 \\
    -a_0 + a_1 - a_2 + 1 &> 0
    \end{aligned}$ & 
    $\begin{aligned}
    |a_0| &< 1 \\
    a_0^2 - a_0 a_2 + a_1 &< 1
    \end{aligned}$ \\
    \hline
\end{ctabular}
\renewcommand{\arraystretch}{1}

\vspace{-0.2cm}


\subsection{Frequenzgang / Bodediagramm / Nyquistdiagramm}{241-243}

Wie bei zeitkontinuierlichen Systemen ergibt sich der \textbf{Frequenzgang}, wenn die UTF $G(z)$ \textbf{entlang der Stabilitätsgrenze ausgewertet} wird.
Die Stabilitätsgrenze ist der Einheitskreis, also alle Stellen $z = e^{\jimg \omega T}$

\vspace{-0.2cm}

$$ G(z) \Big|_{z = e^{\jimg \omega T}} = K \cdot \frac{(z - z_1) (z - z_2) \ldots (z - z_m)}{(z - p_1) (z - p_2) \ldots (z - p_n)} \Bigg|_{z = e^{\jimg \omega T}} $$

Bodediagramm und Nyquistdiagramm können dann \textbf{ähnlich gezeichnet} werden wie bei zeitkontinuierlichen Systemen. 
Es gelten allerdings folgende \textbf{Besonderheiten}:

\smallskip

\begin{outline}
    \1 Darstellungsbereich der Diagramme eingeschränkt auf den Bereich $\omega \in \left[0, \, \omega_N \right]$
        \2 Entspricht Frequenzgang entlang oberer Hälfte des Einheitskreises
        \2 Untere Hälfte ist redundant
    \1 Abtastzeit passend wählen
        \2 Vor allem zu grosse Abtastzeiten führen zu grossen Abweichungen
    \1 Im Bodediagramm existieren \textbf{keine Geradenapproximationen}
    \1 Die Beuteilung der Stabilität wird am offenen Regelkreis durchgeführt
\end{outline}


\subsubsection{Darstellungsbereich}

\begin{minipage}[c]{0.25\columnwidth}
    $$ \omega_N = \frac{\omega_s}{2} $$
\end{minipage}
\hfill
\begin{minipage}[c]{0.7\columnwidth}
    \begin{tabular}{ll|ll}
        $\omega_N$  & Nyquistfrequenz   & $\omega_s$  & Samplingfrequenz  \\
    \end{tabular}
\end{minipage}


